\documentclass[12pt,a4paper]{article}

% Preamble
%\usepackage[utf-8]{inputenc}
\usepackage[margin=1in]{geometry}
\usepackage{amsmath}
\usepackage{amssymb}
\usepackage{hyperref}
\usepackage{xcolor}
\usepackage{booktabs}
\usepackage{array}
\usepackage{setspace}

% Document Metadata
\title{\textbf{GED Math Answer Guide:\\Number Sense and Operations}\\[0.5cm]{\large Detailed Explanations}}
\author{}
\date{}

% Custom Styling
\hypersetup{
    colorlinks=true,
    linkcolor=blue,
    urlcolor=blue
}

\onehalfspacing

\begin{document}

% Title Page
\maketitle
\thispagestyle{empty}
\newpage

% Table of Contents
\tableofcontents
\newpage

% ============================================================================
\section{Basic Level (Questions 1-10)}
% ============================================================================

\subsection*{Question 1: Order of Operations}

\textbf{Correct Answer:} $21$

\subsubsection*{Problem}
Calculate: $8 + 3 \times 5 - 2 = \,?$

\subsubsection*{Step-by-Step Solution}
\begin{enumerate}
    \item Apply the order of operations (PEMDAS/BODMAS)
    \item Multiplication first: $3 \times 5 = 15$
    \item Rewrite the expression: $8 + 15 - 2$
    \item Addition and subtraction (left to right): $8 + 15 = 23$
    \item Final subtraction: $23 - 2 = 21$
\end{enumerate}

\subsubsection*{Mathematical Concepts}
Order of operations ensures consistent calculation results. PEMDAS stands for Parentheses, Exponents, Multiplication and Division (left to right), Addition and Subtraction (left to right).

\subsubsection*{Common Mistakes}
\begin{itemize}
    \item Adding $8 + 3$ first, then multiplying by $5$
    \item Subtracting before adding
    \item Ignoring the order of operations entirely
\end{itemize}

\subsubsection*{How to Avoid}
Memorize PEMDAS and always check for multiplication/division before addition/subtraction.

\subsubsection*{Alternative Method/Verification}
Use a calculator to verify: $8 + 3 \times 5 - 2 = 21$ 

\subsubsection*{Real-World Context}
Budgeting: If you have \$8, earn \$3 per hour for 5 hours, then spend \$2, you have \$21 total.

\vspace{0.5cm}

\subsection*{Question 2: Adding Fractions}

\textbf{Correct Answer:} $\frac{7}{12}$

\subsubsection*{Problem}
Calculate: $\frac{1}{4} + \frac{1}{3} = \,?$

\subsubsection*{Step-by-Step Solution}
\begin{enumerate}
    \item Find the least common denominator (LCD) of 4 and 3
    \item LCD = 12
    \item Convert $\frac{1}{4}$ to twelfths: $\frac{1}{4} = \frac{3}{12}$
    \item Convert $\frac{1}{3}$ to twelfths: $\frac{1}{3} = \frac{4}{12}$
    \item Add the fractions: $\frac{3}{12} + \frac{4}{12} = \frac{7}{12}$
\end{enumerate}

\subsubsection*{Mathematical Concepts}
When adding fractions, denominators must be the same. The LCD is the smallest number divisible by both denominators.

\subsubsection*{Common Mistakes}
\begin{itemize}
    \item Adding numerators and denominators separately: $\frac{1+1}{4+3} = \frac{2}{7}$ (incorrect)
    \item Using 7 as the common denominator instead of 12
\end{itemize}

\subsubsection*{How to Avoid}
Always find the LCD first. Verify that both fractions have the same denominator before adding.

\subsubsection*{Alternative Method/Verification}
Convert to decimals: $0.25 + 0.333... \approx 0.583$, and $\frac{7}{12} \approx 0.583$ 

\subsubsection*{Real-World Context}
Cooking: Using $\frac{1}{4}$ cup of flour and $\frac{1}{3}$ cup of sugar gives $\frac{7}{12}$ cup total ingredients.

\vspace{0.5cm}

\subsection*{Question 3: Subtracting Decimals}

\textbf{Correct Answer:} $3.47$

\subsubsection*{Problem}
Calculate: $8.92 - 5.45 = \,?$

\subsubsection*{Step-by-Step Solution}
\begin{enumerate}
    \item Align decimal points:
    \[\begin{array}{r}
        8.92 \\
        -5.45 \\
        \hline
    \end{array}\]
    \item Subtract from right to left: $2 - 5$ requires borrowing
    \item Borrow 1 from the tenths place: $12 - 5 = 7$
    \item Tenths place: $8 - 4 = 4$
    \item Ones place: $8 - 5 = 3$
    \item Result: $3.47$
\end{enumerate}

\subsubsection*{Mathematical Concepts}
Decimal subtraction follows the same rules as whole number subtraction. Alignment of decimal points is critical for correct computation.

\subsubsection*{Common Mistakes}
\begin{itemize}
    \item Misaligning decimal points
    \item Forgetting to borrow when needed
    \item Subtracting the larger digit from the smaller
\end{itemize}

\subsubsection*{How to Avoid}
Always line up decimal points vertically. Use borrowing notation clearly.

\subsubsection*{Alternative Method/Verification}
Check by addition: $3.47 + 5.45 = 8.92$ 

\subsubsection*{Real-World Context}
Money: If you spend \$5.45 from \$8.92, you have \$3.47 remaining.

\vspace{0.5cm}

\subsection*{Question 4: Converting Fractions to Decimals}

\textbf{Correct Answer:} $0.375$

\subsubsection*{Problem}
Convert $\frac{3}{8}$ to a decimal.

\subsubsection*{Step-by-Step Solution}
\begin{enumerate}
    \item Divide the numerator by the denominator: $3 \div 8$
    \item Long division:
    \[3.000 \div 8 = 0.375\]
    \item $8$ goes into $30$ three times ($8 \times 3 = 24$), remainder $6$
    \item Bring down $0$: $60 \div 8 = 7$ ($8 \times 7 = 56$), remainder $4$
    \item Bring down $0$: $40 \div 8 = 5$ ($8 \times 5 = 40$), remainder $0$
    \item Result: $0.375$
\end{enumerate}

\subsubsection*{Mathematical Concepts}
Any fraction can be converted to a decimal by dividing the numerator by the denominator. Some decimals terminate (like $0.375$), while others repeat.

\subsubsection*{Common Mistakes}
\begin{itemize}
    \item Dividing the denominator by the numerator (incorrect direction)
    \item Stopping division too early
    \item Misplacing the decimal point
\end{itemize}

\subsubsection*{How to Avoid}
Remember: fraction bar means division. Numerator goes inside the division bracket, denominator outside.

\subsubsection*{Alternative Method/Verification}
Multiply both numerator and denominator to get a power of 10 in the denominator:
\[\frac{3}{8} = \frac{3 \times 125}{8 \times 125} = \frac{375}{1000} = 0.375\]

\subsubsection*{Real-World Context}
Measurements: $\frac{3}{8}$ inch equals $0.375$ inches, useful for converting between fractional and decimal measurements.

\vspace{0.5cm}

\subsection*{Question 5: Finding Percentages}

\textbf{Correct Answer:} $25\%$

\subsubsection*{Problem}
15 is what percent of 60?

\subsubsection*{Step-by-Step Solution}
\begin{enumerate}
    \item Set up the percentage formula:
    \[\text{Percentage} = \frac{\text{Part}}{\text{Whole}} \times 100\%\]
    \item Substitute values:
    \[\text{Percentage} = \frac{15}{60} \times 100\%\]
    \item Simplify the fraction: $\frac{15}{60} = \frac{1}{4}$
    \item Calculate: $\frac{1}{4} \times 100\% = 25\%$
\end{enumerate}

\subsubsection*{Mathematical Concepts}
A percentage represents a portion out of 100. The formula is: $\frac{\text{Part}}{\text{Whole}} \times 100\% = \text{Percentage}$

\subsubsection*{Common Mistakes}
\begin{itemize}
    \item Dividing 60 by 15 instead of 15 by 60
    \item Forgetting to multiply by 100
    \item Misidentifying the part and the whole
\end{itemize}

\subsubsection*{How to Avoid}
Remember: the part (smaller number) goes on top, the whole (larger number) on the bottom.

\subsubsection*{Alternative Method/Verification}
Use proportion: $\frac{15}{60} = \frac{x}{100}$, so $x = 25$ 

\subsubsection*{Real-World Context}
Sales: 15 out of 60 items sold represents a 25\% sales conversion rate.

\vspace{0.5cm}

\subsection*{Question 6: Multiplying Fractions}

\textbf{Correct Answer:} $\frac{2}{15}$

\subsubsection*{Problem}
Calculate: $\frac{2}{3} \times \frac{1}{5} = \,?$

\subsubsection*{Step-by-Step Solution}
\begin{enumerate}
    \item Multiply numerators: $2 \times 1 = 2$
    \item Multiply denominators: $3 \times 5 = 15$
    \item Result: $\frac{2}{15}$
    \item Check if simplifiable: GCD(2, 15) = 1, so $\frac{2}{15}$ is in simplest form
\end{enumerate}

\subsubsection*{Mathematical Concepts}
To multiply fractions, multiply numerators together and denominators together. Always simplify to lowest terms.

\subsubsection*{Common Mistakes}
\begin{itemize}
    \item Adding instead of multiplying
    \item Finding a common denominator (only needed for addition/subtraction)
    \item Not simplifying the final answer
\end{itemize}

\subsubsection*{How to Avoid}
Use the rule: "multiply straight across." Check if common factors exist before multiplying to simplify early.

\subsubsection*{Alternative Method/Verification}
Convert to decimals: $0.667 \times 0.2 = 0.133...$, and $\frac{2}{15} = 0.133...$ 

\subsubsection*{Real-World Context}
Recipes: $\frac{2}{3}$ of $\frac{1}{5}$ cup of salt equals $\frac{2}{15}$ cup.

\vspace{0.5cm}

\subsection*{Question 7: Dividing Fractions}

\textbf{Correct Answer:} $\frac{8}{9}$

\subsubsection*{Problem}
Calculate: $\frac{2}{3} \div \frac{3}{4} = \,?$

\subsubsection*{Step-by-Step Solution}
\begin{enumerate}
    \item Write the division problem: $\frac{2}{3} \div \frac{3}{4}$
    \item Flip (invert) the second fraction: $\frac{3}{4}$ becomes $\frac{4}{3}$
    \item Change division to multiplication: $\frac{2}{3} \times \frac{4}{3}$
    \item Multiply numerators: $2 \times 4 = 8$
    \item Multiply denominators: $3 \times 3 = 9$
    \item Result: $\frac{8}{9}$
\end{enumerate}

\subsubsection*{Mathematical Concepts}
Dividing by a fraction is the same as multiplying by its reciprocal. The reciprocal of $\frac{a}{b}$ is $\frac{b}{a}$.

\subsubsection*{Common Mistakes}
\begin{itemize}
    \item Flipping the first fraction instead of the second
    \item Forgetting to invert the divisor entirely
    \item Adding or subtracting instead of multiplying after inversion
\end{itemize}

\subsubsection*{How to Avoid}
Use the phrase: ``Keep, Change, Flip''---keep the first fraction, change $\div$ to $\times$, flip the second fraction.

\subsubsection*{Alternative Method/Verification}
Multiply by reciprocal: $\frac{2}{3} \times \frac{4}{3} = \frac{8}{9}$ 

\subsubsection*{Real-World Context}
Sharing: If you have $\frac{2}{3}$ of a pizza and divide it into portions of $\frac{3}{4}$, you get $\frac{8}{9}$ portions.

\vspace{0.5cm}

\subsection*{Question 8: Exponents}

\textbf{Correct Answer:} $32$

\subsubsection*{Problem}
Calculate: $2^5 = \,?$

\subsubsection*{Step-by-Step Solution}
\begin{enumerate}
    \item Understand exponent notation: $2^5$ means $2$ multiplied by itself $5$ times
    \item Write out the multiplication: $2 \times 2 \times 2 \times 2 \times 2$
    \item Calculate step by step:
    \begin{align*}
        2 \times 2 &= 4 \\
        4 \times 2 &= 8 \\
        8 \times 2 &= 16 \\
        16 \times 2 &= 32
    \end{align*}
    \item Result: $32$
\end{enumerate}

\subsubsection*{Mathematical Concepts}
An exponent (or power) tells how many times the base is multiplied by itself. In $2^5$, the base is $2$ and the exponent is $5$.

\subsubsection*{Common Mistakes}
\begin{itemize}
    \item Multiplying the base by the exponent: $2 \times 5 = 10$ (incorrect)
    \item Computing only partial repetitions
    \item Confusing exponents with coefficients
\end{itemize}

\subsubsection*{How to Avoid}
Expand the expression fully and multiply step by step. Use a calculator for verification.

\subsubsection*{Alternative Method/Verification}
Use a calculator: $2^5 = 32$ 

\subsubsection*{Real-World Context}
Data storage: A computer address with $2^5$ possible values has 32 unique combinations.

\vspace{0.5cm}

\subsection*{Question 9: Square Roots}

\textbf{Correct Answer:} $8$

\subsubsection*{Problem}
Calculate: $\sqrt{64} = \,?$

\subsubsection*{Step-by-Step Solution}
\begin{enumerate}
    \item Understand square root: Find the number that, when multiplied by itself, equals 64
    \item Think: What times what equals 64?
    \item Test: $8 \times 8 = 64$
    \item Result: $\sqrt{64} = 8$
\end{enumerate}

\subsubsection*{Mathematical Concepts}
A square root is the inverse operation of squaring. If $a^2 = b$, then $\sqrt{b} = a$.

\subsubsection*{Common Mistakes}
\begin{itemize}
    \item Dividing by 2: $64 \div 2 = 32$ (incorrect)
    \item Confusing square root with division by 2
    \item Getting negative square roots (though $-8$ is also technically valid)
\end{itemize}

\subsubsection*{How to Avoid}
Learn perfect squares: $1, 4, 9, 16, 25, 36, 49, 64, 81, 100, \ldots$

\subsubsection*{Alternative Method/Verification}
Check by squaring: $8^2 = 8 \times 8 = 64$ 

\subsubsection*{Real-World Context}
Geometry: A square with area $64$ square feet has sides of length $8$ feet.

\vspace{0.5cm}

\subsection*{Question 10: Absolute Value}

\textbf{Correct Answer:} $5$

\subsubsection*{Problem}
Calculate: $|-5| = \,?$

\subsubsection*{Step-by-Step Solution}
\begin{enumerate}
    \item Understand absolute value: It is the distance from zero on a number line, always positive
    \item The number $-5$ is 5 units away from zero
    \item Result: $|-5| = 5$
\end{enumerate}

\subsubsection*{Mathematical Concepts}
Absolute value $|x|$ is always non-negative. It represents the magnitude or distance from zero, regardless of direction.

\subsubsection*{Common Mistakes}
\begin{itemize}
    \item Returning the negative value: $|-5| = -5$ (incorrect)
    \item Confusing absolute value with the number itself
\end{itemize}

\subsubsection*{How to Avoid}
Remember: absolute value makes all numbers positive (or zero).

\subsubsection*{Alternative Method/Verification}
Check: $|-5| = |(-1) \times 5| = 1 \times |5| = 5$ 

\subsubsection*{Real-World Context}
Temperature: An absolute value of $-5°$ represents 5 degrees below zero.

\newpage

% ============================================================================
\section{Intermediate Level (Questions 11-20)}
% ============================================================================

\subsection*{Question 11: Complex Order of Operations}

\textbf{Correct Answer:} $2$

\subsubsection*{Problem}
Calculate: $24 \div (4 + 2) - 3 + 1 = \,?$

\subsubsection*{Step-by-Step Solution}
\begin{enumerate}
    \item Parentheses first: $4 + 2 = 6$
    \item Rewrite: $24 \div 6 - 3 + 1$
    \item Division: $24 \div 6 = 4$
    \item Rewrite: $4 - 3 + 1$
    \item Subtract: $4 - 3 = 1$
    \item Add: $1 + 1 = 2$
\end{enumerate}

\subsubsection*{Mathematical Concepts}
PEMDAS/BODMAS: Parentheses/Brackets first, then Exponents/Orders, then Multiplication and Division (left to right), then Addition and Subtraction (left to right).

\subsubsection*{Common Mistakes}
\begin{itemize}
    \item Ignoring parentheses and dividing 24 by 4 first
    \item Performing operations in incorrect sequence
    \item Subtracting before adding
\end{itemize}

\subsubsection*{How to Avoid}
Always start with parentheses. Work left to right for operations at the same level.

\subsubsection*{Alternative Method/Verification}
Use a calculator: $24 \div (4 + 2) - 3 + 1 = 2$ 

\subsubsection*{Real-World Context}
Budget calculation: Sharing \$24 among 6 people (\$4 + 2 people), then subtracting a \$3 fee and adding \$1 bonus gives \$2 per person.

\vspace{0.5cm}

\subsection*{Question 12: Converting Percentages}

\textbf{Correct Answer:} $\$37.50$

\subsubsection*{Problem}
A shirt costs \$50 with a 25\% discount. What is the final price?

\subsubsection*{Step-by-Step Solution}
\begin{enumerate}
    \item Calculate the discount amount: $25\% \text{ of } \$50 = 0.25 \times 50 = 12.50$
    \item Subtract from original price: $50 - 12.50 = 37.50$
    \item Final price: $\$37.50$
\end{enumerate}

\subsubsection*{Mathematical Concepts}
A percentage discount reduces the original price. Formula: Final Price = Original Price $-$ (Percentage $\times$ Original Price)

\subsubsection*{Common Mistakes}
\begin{itemize}
    \item Calculating 25\% of 25 instead of 25\% of 50
    \item Forgetting to subtract the discount from the original
    \item Using 250\% instead of 0.25
\end{itemize}

\subsubsection*{How to Avoid}
Convert percentage to decimal (divide by 100), multiply by the original amount, then subtract.

\subsubsection*{Alternative Method/Verification}
If 25\% off, you pay 75\%: $0.75 \times 50 = 37.50$ 

\subsubsection*{Real-World Context}
Shopping: A \$50 shirt with a 25\% sale saves you \$12.50.

\vspace{0.5cm}

\subsection*{Question 13: Ratio and Proportion}

\textbf{Correct Answer:} $12$ cups flour

\subsubsection*{Problem}
A recipe calls for flour to sugar in the ratio $3:2$. If you use $8$ cups of sugar, how much flour do you need?

\subsubsection*{Step-by-Step Solution}
\begin{enumerate}
    \item Set up the proportion: $\frac{3 \text{ flour}}{2 \text{ sugar}} = \frac{x \text{ flour}}{8 \text{ sugar}}$
    \item Cross-multiply: $3 \times 8 = 2 \times x$
    \item Simplify: $24 = 2x$
    \item Solve for $x$: $x = \frac{24}{2} = 12$
    \item Answer: $12$ cups flour
\end{enumerate}

\subsubsection*{Mathematical Concepts}
Ratios compare quantities. A proportion states that two ratios are equal. Cross-multiplication solves proportions.

\subsubsection*{Common Mistakes}
\begin{itemize}
    \item Inverting the ratio parts
    \item Not cross-multiplying correctly
    \item Confusing the numerator and denominator placement
\end{itemize}

\subsubsection*{How to Avoid}
Set up the proportion clearly with consistent placement. Cross-multiply: $\frac{a}{b} = \frac{c}{d}$ gives $ad = bc$.

\subsubsection*{Alternative Method/Verification}
Scale factor: $8 \div 2 = 4$, so flour is $3 \times 4 = 12$ 

\subsubsection*{Real-World Context}
Cooking: Scaling recipes using ratios ensures ingredients maintain proper proportions.

\vspace{0.5cm}

\subsection*{Question 14: Multi-Step Percentage}

\textbf{Correct Answer:} $\$216,000$

\subsubsection*{Problem}
A house costs \$200,000. After a 20\% increase, then a 10\% decrease, what is the final price?

\subsubsection*{Step-by-Step Solution}
\begin{enumerate}
    \item Calculate 20\% increase: $20\% \text{ of } 200,000 = 0.20 \times 200,000 = 40,000$
    \item New price after increase: $200,000 + 40,000 = 240,000$
    \item Calculate 10\% decrease: $10\% \text{ of } 240,000 = 0.10 \times 240,000 = 24,000$
    \item Final price: $240,000 - 24,000 = 216,000$
    \item Answer: $\$216,000$
\end{enumerate}

\subsubsection*{Mathematical Concepts}
When applying multiple percentage changes, apply them sequentially. The second percentage applies to the new amount, not the original.

\subsubsection*{Common Mistakes}
\begin{itemize}
    \item Adding 20\% and subtracting 10\% from the original (ignoring compounding)
    \item Calculating the 10\% decrease from the original price
    \item Arithmetic errors in intermediate steps
\end{itemize}

\subsubsection*{How to Avoid}
Work step by step. Always use the current amount as the base for the next percentage calculation.

\subsubsection*{Alternative Method/Verification}
Multiply by factors: $200,000 \times 1.20 \times 0.90 = 216,000$ 

\subsubsection*{Real-World Context}
Real estate: A house appreciates 20\% then depreciates 10\%, not returning to the original value.

\vspace{0.5cm}

\subsection*{Question 15: Fraction to Percent}

\textbf{Correct Answer:} $60\%$

\subsubsection*{Problem}
Convert $\frac{3}{5}$ to a percent.

\subsubsection*{Step-by-Step Solution}
\begin{enumerate}
    \item Use the formula: Percentage $= \frac{\text{Fraction}}{1} \times 100\%$
    \item Substitute: $\frac{3}{5} \times 100\%$
    \item Simplify: $\frac{3 \times 100}{5}\% = \frac{300}{5}\% = 60\%$
\end{enumerate}

\subsubsection*{Mathematical Concepts}
A fraction can be converted to a percentage by multiplying by 100. Alternatively, find an equivalent fraction with denominator 100.

\subsubsection*{Common Mistakes}
\begin{itemize}
    \item Forgetting to multiply by 100
    \item Dividing 100 by 5 and 3 separately
    \item Confusing the operation direction
\end{itemize}

\subsubsection*{How to Avoid}
Always multiply the fraction by 100 to convert to percent.

\subsubsection*{Alternative Method/Verification}
Equivalent fraction: $\frac{3}{5} = \frac{60}{100} = 60\%$ 

\subsubsection*{Real-World Context}
Grades: Scoring $\frac{3}{5}$ questions correctly equals $60\%$.

\vspace{0.5cm}

\subsection*{Question 16: Negative Numbers}

\textbf{Correct Answer:} $-8$

\subsubsection*{Problem}
Calculate: $-3 - 5 = \,?$

\subsubsection*{Step-by-Step Solution}
\begin{enumerate}
    \item Interpret subtraction as adding a negative: $-3 - 5 = -3 + (-5)$
    \item Both numbers are negative; add their absolute values: $3 + 5 = 8$
    \item Result is negative: $-8$
    \item Answer: $-8$
\end{enumerate}

\subsubsection*{Mathematical Concepts}
Subtracting a positive number is the same as adding a negative. When adding two negatives, the result is negative with a magnitude equal to the sum of the absolute values.

\subsubsection*{Common Mistakes}
\begin{itemize}
    \item Getting $-2$ by subtracting 5 from 3 (ignoring the initial negative)
    \item Incorrect sign handling in intermediate steps
\end{itemize}

\subsubsection*{How to Avoid}
Think of the number line: starting at $-3$, move 5 units left to reach $-8$.

\subsubsection*{Alternative Method/Verification}
Use a number line or calculator: $-3 - 5 = -8$ 

\subsubsection*{Real-World Context}
Temperature: If it's $-3°$ and drops 5 more degrees, it becomes $-8°$.

\vspace{0.5cm}

\subsection*{Question 17: Multiplying Negatives}

\textbf{Correct Answer:} $-12$

\subsubsection*{Problem}
Calculate: $3 \times (-4) = \,?$

\subsubsection*{Step-by-Step Solution}
\begin{enumerate}
    \item Recall the rule: positive times negative equals negative
    \item Multiply absolute values: $3 \times 4 = 12$
    \item Apply the sign: negative
    \item Answer: $-12$
\end{enumerate}

\subsubsection*{Mathematical Concepts}
Sign rules for multiplication:
\begin{itemize}
    \item Positive $\times$ Positive = Positive
    \item Positive $\times$ Negative = Negative
    \item Negative $\times$ Negative = Positive
\end{itemize}

\subsubsection*{Common Mistakes}
\begin{itemize}
    \item Treating the negative as a positive: $3 \times 4 = 12$
    \item Using an incorrect sign rule
\end{itemize}

\subsubsection*{How to Avoid}
Remember: unlike signs give a negative product. Same signs give a positive product.

\subsubsection*{Alternative Method/Verification}
Repeated addition: $3 \times (-4) = (-4) + (-4) + (-4) = -12$ 

\subsubsection*{Real-World Context}
Debt: Having 3 debts of \$4 each equals a total debt of \$12 (negative balance).

\vspace{0.5cm}

\subsection*{Question 18: Multiple Fractions}

\textbf{Correct Answer:} $\frac{2}{3}$

\subsubsection*{Problem}
Calculate: $\frac{1}{2} + \frac{1}{6} = \,?$

\subsubsection*{Step-by-Step Solution}
\begin{enumerate}
    \item Find LCD of 2 and 6: LCD = 6
    \item Convert $\frac{1}{2}$ to sixths: $\frac{1}{2} = \frac{3}{6}$
    \item Keep $\frac{1}{6}$ as is
    \item Add: $\frac{3}{6} + \frac{1}{6} = \frac{4}{6}$
    \item Simplify: $\frac{4}{6} = \frac{2}{3}$
\end{enumerate}

\subsubsection*{Mathematical Concepts}
LCD is the smallest number divisible by all denominators. Simplification reduces fractions to lowest terms by dividing by the GCD.

\subsubsection*{Common Mistakes}
\begin{itemize}
    \item Using 12 as the LCD instead of 6
    \item Forgetting to simplify the final answer
    \item Adding denominators
\end{itemize}

\subsubsection*{How to Avoid}
Find the LCD by listing multiples. Always simplify the final fraction.

\subsubsection*{Alternative Method/Verification}
Decimals: $0.5 + 0.167 = 0.667$, and $\frac{2}{3} = 0.667$ 

\subsubsection*{Real-World Context}
Portions: $\frac{1}{2}$ pizza plus $\frac{1}{6}$ pizza equals $\frac{2}{3}$ pizza.

\vspace{0.5cm}

\subsection*{Question 19: Tax Calculation}

\textbf{Correct Answer:} $\$65.52$

\subsubsection*{Problem}
A book costs \$58.50. With a 12\% sales tax, what is the total cost?

\subsubsection*{Step-by-Step Solution}
\begin{enumerate}
    \item Calculate tax amount: $12\% \text{ of } 58.50 = 0.12 \times 58.50 = 7.02$
    \item Add tax to original price: $58.50 + 7.02 = 65.52$
    \item Total: $\$65.52$
\end{enumerate}

\subsubsection*{Mathematical Concepts}
Sales tax is a percentage added to the original price. Formula: Total = Original Price + (Tax Rate $\times$ Original Price)

\subsubsection*{Common Mistakes}
\begin{itemize}
    \item Calculating 12\% of 12 instead of 12\% of 58.50
    \item Forgetting to add the tax to the original price
    \item Rounding prematurely
\end{itemize}

\subsubsection*{How to Avoid}
Convert the percentage to decimal, multiply by the original amount, then add to find the total.

\subsubsection*{Alternative Method/Verification}
Total is 112\% of original: $1.12 \times 58.50 = 65.52$ 

\subsubsection*{Real-World Context}
Shopping: Most purchases include sales tax, increasing the final price.

\vspace{0.5cm}

\subsection*{Question 20: Comparing Forms}

\textbf{Correct Answer:} $\frac{3}{4} = 0.75 = 75\%$

\subsubsection*{Problem}
Show that $\frac{3}{4}$, $0.75$, and $75\%$ are equivalent.

\subsubsection*{Step-by-Step Solution}
\begin{enumerate}
    \item Convert fraction to decimal: $\frac{3}{4} = 3 \div 4 = 0.75$
    \item Convert decimal to percent: $0.75 \times 100\% = 75\%$
    \item All three forms represent the same value
\end{enumerate}

\subsubsection*{Mathematical Concepts}
Fractions, decimals, and percentages are interchangeable representations of the same quantity.

\subsubsection*{Common Mistakes}
\begin{itemize}
    \item Converting fractions incorrectly to decimals
    \item Confusing decimal multiplication when converting to percent
    \item Not recognizing equivalent forms
\end{itemize}

\subsubsection*{How to Avoid}
Practice converting between all three forms systematically.

\subsubsection*{Comparison Table}

\begin{center}
\begin{tabular}{|c|c|c|}
\toprule
\textbf{Fraction} & \textbf{Decimal} & \textbf{Percent} \\
\midrule
$\frac{1}{4}$ & $0.25$ & $25\%$ \\
$\frac{1}{2}$ & $0.50$ & $50\%$ \\
$\frac{3}{4}$ & $0.75$ & $75\%$ \\
$\frac{1}{5}$ & $0.20$ & $20\%$ \\
$\frac{2}{5}$ & $0.40$ & $40\%$ \\
$\frac{3}{5}$ & $0.60$ & $60\%$ \\
\bottomrule
\end{tabular}
\end{center}

\subsubsection*{Real-World Context}
Understanding different representations helps in interpreting data presented in various formats.

\newpage

% ============================================================================
\section{Advanced Level (Questions 21-30)}
% ============================================================================

\subsection*{Question 21: Complex Fractions}

\textbf{Correct Answer:} $\frac{19}{24}$

\subsubsection*{Problem}
Calculate: $\frac{5}{8} + \frac{1}{6} = \,?$

\subsubsection*{Step-by-Step Solution}
\begin{enumerate}
    \item Find LCD of 8 and 6:
    \begin{align*}
        \text{Multiples of } 8 &: 8, 16, 24, 32, \ldots \\
        \text{Multiples of } 6 &: 6, 12, 18, 24, \ldots \\
        \text{LCD} &= 24
    \end{align*}
    \item Convert $\frac{5}{8}$ to twenty-fourths: $\frac{5}{8} = \frac{5 \times 3}{8 \times 3} = \frac{15}{24}$
    \item Convert $\frac{1}{6}$ to twenty-fourths: $\frac{1}{6} = \frac{1 \times 4}{6 \times 4} = \frac{4}{24}$
    \item Add: $\frac{15}{24} + \frac{4}{24} = \frac{19}{24}$
    \item Check if simplifiable: GCD(19, 24) = 1, so $\frac{19}{24}$ is in simplest form
\end{enumerate}

\subsubsection*{Mathematical Concepts}
Finding the LCD requires identifying the least common multiple of the denominators. Prime factorization can help find the LCD systematically.

\subsubsection*{Common Mistakes}
\begin{itemize}
    \item Using 48 as the LCD instead of 24
    \item Incorrect conversion of fractions
    \item Arithmetic errors in finding the LCD
\end{itemize}

\subsubsection*{How to Avoid}
List multiples of each denominator until you find a common one. Use prime factorization for complex cases.

\subsubsection*{Alternative Method/Verification}
Decimals: $0.625 + 0.167 = 0.792$, and $\frac{19}{24} \approx 0.792$ 

\subsubsection*{Real-World Context}
Measurement: Combining $\frac{5}{8}$ inch and $\frac{1}{6}$ inch yields $\frac{19}{24}$ inch.

\vspace{0.5cm}

\subsection*{Question 22: Exponents}

\textbf{Correct Answer:} $41$

\subsubsection*{Problem}
Calculate: $5^2 + 4^2 = \,?$

\subsubsection*{Step-by-Step Solution}
\begin{enumerate}
    \item Evaluate $5^2$: $5 \times 5 = 25$
    \item Evaluate $4^2$: $4 \times 4 = 16$
    \item Add the results: $25 + 16 = 41$
\end{enumerate}

\subsubsection*{Mathematical Concepts}
Exponentiation is evaluated before addition. This relates to the Pythagorean theorem: $a^2 + b^2 = c^2$.

\subsubsection*{Common Mistakes}
\begin{itemize}
    \item Computing $5^2 + 4^2 = (5 + 4)^2 = 81$ (incorrect grouping)
    \item Multiplying instead of squaring
    \item Adding first, then squaring: $(5 + 4)^2 = 81$
\end{itemize}

\subsubsection*{How to Avoid}
Evaluate exponents individually before performing other operations.

\subsubsection*{Alternative Method/Verification}
Calculator: $5^2 + 4^2 = 25 + 16 = 41$ 

\subsubsection*{Real-World Context}
Pythagorean theorem: A right triangle with legs of 5 and 4 units has a hypotenuse of $\sqrt{41}$ units.

\vspace{0.5cm}

\subsection*{Question 23: Square Roots}

\textbf{Correct Answer:} $13$ ft

\subsubsection*{Problem}
A room has an area of $169$ square feet. If the room is square-shaped, what is the length of one side?

\subsubsection*{Step-by-Step Solution}
\begin{enumerate}
    \item For a square, Area $= \text{side}^2$
    \item Set up the equation: $\text{side}^2 = 169$
    \item Take the square root: $\text{side} = \sqrt{169}$
    \item Calculate: $\sqrt{169} = 13$ (since $13 \times 13 = 169$)
    \item Answer: $13$ feet
\end{enumerate}

\subsubsection*{Mathematical Concepts}
The square root is the inverse of squaring. For a square, the side length is the square root of the area.

\subsubsection*{Common Mistakes}
\begin{itemize}
    \item Dividing the area by 2: $169 \div 2 = 84.5$ (incorrect)
    \item Using $\frac{169}{4}$ as the side length
    \item Not recognizing 13 as the square root of 169
\end{itemize}

\subsubsection*{How to Avoid}
Memorize perfect squares up to $20^2 = 400$.

\subsubsection*{Alternative Method/Verification}
Check: $13 \times 13 = 169$ 

\subsubsection*{Real-World Context}
Flooring: Determining room dimensions from area measurements requires taking square roots.

\vspace{0.5cm}

\subsection*{Question 24: Compound Percent}

\textbf{Correct Answer:} $\$115.56$

\subsubsection*{Problem}
An investment of \$100 earns 8\% interest, then the new amount earns 7\% interest. What is the final amount?

\subsubsection*{Step-by-Step Solution}
\begin{enumerate}
    \item Initial investment: \$100
    \item After first year (8\% interest): $100 \times 1.08 = 108$
    \item After second year (7\% on new amount): $108 \times 1.07 = 115.56$
    \item Final amount: \$115.56
\end{enumerate}

\subsubsection*{Mathematical Concepts}
Compound interest applies interest to the accumulated amount, not the original. Formula: Final Amount $= P(1 + r_1)(1 + r_2)$

\subsubsection*{Common Mistakes}
\begin{itemize}
    \item Adding interest rates: $8\% + 7\% = 15\%$, then applying $15\%$ to original
    \item Calculating the second interest on the original amount
    \item Rounding intermediate steps
\end{itemize}

\subsubsection*{How to Avoid}
Apply each interest rate sequentially to the current amount.

\subsubsection*{Alternative Method/Verification}
$100 \times 1.08 \times 1.07 = 100 \times 1.1556 = 115.56$ 

\subsubsection*{Real-World Context}
Banking: Compound interest is how savings accounts and investments grow over multiple periods.

\vspace{1cm}

\noindent
\textit{This completes the GED Math Answer Guide for Number Sense and Operations.}

\end{document}